% 五、课题执行计划。
% 拟于2027年1月至2028年12月实施；具体日期依正式批复。
\subsection{年度研究计划}
本课题拟于2027年1月至2028年12月实施，共24个月，具体日期依批复执行。优先完成MineNetCD核心实验、结果复核及论文写作，再视余量开展外部验证和额外方法复现。

\textbf{第一年度：完成主数据、失效诊断和核心机制验证。}
第1—3个月验收归档与矿区来源，审计划分和权重覆盖，固定50/10/10/30名单、特征层及掩码投影。复现ResNet-18孪生与DonaNet配置并计时，测试保持封存。完成至多12次短训练，形成误差清单，登记训练集$D_s$与小变化面积界点；仅用训练、验证数据预查端点样本和强度匹配规则。来源或标签有疑问时先处理质量问题。

第4—9个月完成四组约束的12次主训练和DonaNet两组的6次近邻训练，开展等面积、平滑及固定模型梯度诊断。最多6次补充训练优先用于强度匹配，额外模型和外部训练后置且不得突破总额度。用验证预测冻结阈值敏感性规则，保存逐矿区概率、参数和日志，同步撰写初稿。

第10—12个月固定检测器，完成独立校准与错误排序原型，并复用无约束、全图和背景组模型检查残留错误的构成。核对分母及数据用途，完善初稿，最终结果待后续评估补入。第一年度提交可复现实验、误差分析和核心对照记录，按要求准备中期材料，如实报告无改善及尚未解决的问题。

\textbf{第二年度：完成可靠性评价、统计分析与成果发表。}
第13—15个月按预先确定的方案完成矿区评估，在10\%主复核面积下评价效果，并进行0、5\%、20\%补充分析，完成10000次矿区配对重采样，集中检查高置信度漏检、相对小变化及困难矿区。核心结果足以独立成稿时，争取第15个月完成首投；否则说明还缺少哪些证据，并调整投稿安排。完成主数据核心对照后，再在预留额度内决定是否开展EuroMineNet内部验证及专业人员案例复核；这两项补充工作不能代替主数据上的比较。

第16—18个月汇总误报与召回权衡、失效条件及样本不足情况，处理审稿意见；第19—24个月修改论文并争取按指南要求发表，归档代码、数据划分、环境和统计记录，提交结题报告。无改善及召回损伤结果同样保留并分析；预留审稿修改时间，按管理要求标注成果署名、基金名称和批准号。
